% !PRINTED 309
\weberpage{309}
%% [page furniture, not body text -- commented out 2026-08-10] Page furniture (running
%% head, folio, signature) is excluded from this edition by standing policy.
%% \noindent \S.\,80. Equations of the 8th Degree.

%<<CONT
so that $R$ is the entire ternary linear congruence group for the modulus 2.

We consider a system of eight independent variables $X_{z_1,z_2,z_3}$, in which the indices $z_1, z_2, z_3$ are to be taken with respect to the modulus 2.

From these we form the seven functions (\gesperrt{Lagrange} resolvents, Bd.\,I, \S.\,164).
\weq{1}{\overset{z}{\Sigma} (-1)^{\Sigma z\xi} X_{z_1,z_2,z_3} = \Psi_{\xi_1,\xi_2,\xi_3},}
in which, for abbreviation,
\[
\Sigma z\xi = z_1\xi_1 + z_2\xi_2 + z_3\xi_3
\]
is set, and $\xi_1, \xi_2, \xi_3$ each run through a complete system of residues with respect to the modulus 2, but with the exclusion of the combination $0, 0, 0$.

If we apply the substitutions of the group $Q$ to the indices $z_i$ of the quantities $X$, the $X$ undergo the permutations of a group, which we likewise denote by $Q$. The changes of the $\Psi$ that correspond to these permutations follow from (1), if we apply to the indices of the $X$ the substitution $(z_i, z_i+h)$:
\weq{2}{\Sigma (-1)^{\Sigma z\xi} X_{z_1+h_1, z_2+h_2, z_3+h_3} = \overset{z}{\Sigma} (-1)^{\Sigma(z+h)\xi} X_{z_1,z_2,z_3} = (-1)^{\Sigma h\xi} \Psi_{\xi_1,\xi_2,\xi_3}.}

\gesperrt{1. The functions $\Psi$ therefore change only their sign under the permutations of $Q$, and the squares of the $\Psi$ remain unchanged by $Q$.}

We have further to determine the influence of a substitution $\sigma$ from $S$ upon the functions $\Psi$.

For this purpose let us denote by $\varrho$ the \gesperrt{transposed substitution} to $\sigma$, and set
\weq{3}{z' = \sigma(z), \quad \xi = \varrho(\xi'), \quad \xi' = \varrho^{-1}(\xi),}
then (\S.\,37, 9.):
\weq{4}{\Sigma \xi z = \Sigma \xi' z'.}

If we make the substitutions $\sigma$ of the group $S$ in the indices of $X$, we obtain a permutation group of the $X$, to be denoted by $S$, and from (1) it follows that
\[
\overset{z}{\Sigma} (-1)^{\Sigma z\xi} X_{z_1', z_2', z_3'} = \overset{z}{\Sigma} (-1)^{\Sigma \xi' z'} X_{z_1', z_2', z_3'},
\]
%CONT>>
